\documentclass[11pt]{article}
\usepackage{manual_docs_style}

\begin{document}

\docTitle{Terminology and Common Utilities}
\phantomsection\label{docstart:terminology_common_utilities}

\section*{Terminology}
\begin{itemize}
\item (\termSRD) For values \(a\), \(b\), and \(\varepsilon\), \termSRD is defined as
\[
\operatorname{SRD}(a,b;\varepsilon)
=
\frac{2|a-b|}{|a|+|b|+\varepsilon}.
\]
When \(\varepsilon=0\) and \(a,b\neq 0\), this equals the harmonic mean of the one-sided relative differences
\[
\frac{|a-b|}{|a|}
\quad\text{and}\quad
\frac{|a-b|}{|b|},
\]

\item (\termMinAbsDist): minimum absolute distance/change required between values.
\item (\termNoParameterChangeClass): No intervention class (i.e. no parameter is changed during the simulation). Class indicating that none of the exposed intervention parameters changes during the observed simulation interval. In the text we will use the symbol \termNoParameterChangeClass excepts in prompts.
In this case we use the corresponding agent-facing class is "no intervention". However in the paper we submitted to NeuroNIPS we used "no parameter change"
\item (\termDescLevel): It describes the amount of textual description provided to the agent. Its allowed values are \code{none} and \code{high}.
\item (\termSimulator): One of physical simulator from which we generate data. The allowed values are defined in \code{ALLOWED_TSENV_MODELS}.
\item (\termSimulatorRecipe): One configuration that can be given as input to the simulate stage to generate.
\item (\termRunID): An identifier for a single run. 
\item (\termFamily): A \termFamily is the baseline-centered group of run nodes generated from one sampled baseline recipe. 
It contains the baseline run, all intervention children derived from that baseline, and the matched time-zero/static-change probe for each intervention. 
All nodes in the group share the same \code{family_id}, which is \termRunID of the baseline.
\item (\termTaskMode): The required form of the agent's final answer. 
Its allowed values are \code{direct} and \code{code}. 
In \code{direct} mode, the agent must directly write the predictions for the test samples. In \code{code} mode, the agent must write an interpretable \code{rule.py} file defining a prediction function that is evaluated on the test samples and held-out samples.
\item (\termNoiseLevel): The observation-noise profile applied when materializing a sample.
Its allowed values are \code{none}, \code{low}, and \code{high}, and these values are passed to \code{NOISE_LEVEL_MULTIPLIER} which maps them to intergers that are then multiplied by NOISE_DICT.
Currently: \code{none} -> \code{0}, \code{low} -> \code{1}, and \code{high} -> \code{4}.
In the \code{low} noise regime each signal must have a signal to noise ratio higher than 30dB.
\end{itemize}

\section*{Common Utilities}

\begin{algorithm}[H]
  \caption{\code{materialize(uuid_path, noise_level="none", seed, uuid_reference_path=None) -> (df, noise_analysis)}}
\label{def:materialize-utility}
\textbf{Purpose.} Loads a saved run artifact, optionally applies the requested observation-noise profile, anonymizes columns to \code{col1}, \code{col2}, \ldots, \code{colN}, and returns the agent-visible dataframe.

\textbf{Used by.} Cheap Filter Metrics, Agentic Rollout, and Final Artifact Scoring.
\par\smallskip
\begin{algorithmic}[1]
\State \code{src} $\gets$ read parquet file at \code{uuid_path}
\State \code{ref} $\gets$ \code{None}
\State \code{noise_analysis} $\gets$ \code{None}
\If{\code{uuid_reference_path} is not \code{None}}
    \State \code{ref} $\gets$ read parquet file at \code{uuid_reference_path}
\EndIf
\If{\code{noise_level} $\neq$ \code{"none"}}
    \State \code{(noisy_df, noise_analysis)} $\gets$ \code{add_noise(src, seed, noise_level, ref=ref)}
    \State \code{df} $\gets$ \code{noisy_df}
\Else
    \State \code{df} $\gets$ \code{src}
\EndIf
\State Rename all columns of \code{df} to \code{col1}, \code{col2}, \ldots, \code{colN}
\Assert \code{colN} is time (i.e. regularly sampled)
\State \Return \code{(df, noise_analysis)}
\end{algorithmic}
\end{algorithm}

Where \code{add_noise} is defined for each \termSimulator in \code{noise_adder.py}.

Each \termSimulator's \code{noise_adder.py} defines \code{NOISE_DICT}, whose values give the \termSimulator-specific base magnitudes for the low-noise regime.
\begin{algorithm}[H]
\caption{\code{add_noise(...)}}
\label{def:add-noise-utility}
\textbf{Purpose.} Applies the \termSimulator-specific low or high observation-noise profile to a clean simulator trajectory and optionally to its matched reference trajectory.

\textbf{Used by.} \code{materialize(...)}, Cheap Filter Metrics, Agentic Rollout, and noise-analysis reporting.
\par\smallskip
\begin{algorithmic}[1]
\Assert{\code{noise_level} is in \code{"none","low","high"}}
\State \code{current_seed} $\gets$ \code{seed}
\State \code{scale} $\gets$ \code{NOISE_LEVEL_MULTIPLIER[noise_level]}
\State \code{src_noisy} $\gets$ apply \termSimulator-specific noise to \code{src} using \code{current_seed}, \code{scale}, \code{NOISE_DICT}
\If{\code{ref} is not \code{None}}
    \State \code{ref_noisy} $\gets$ apply \termSimulator-specific noise to \code{ref} using \code{current_seed+10}, \code{scale}, \code{NOISE_DICT}
\Else
    \State \code{ref_noisy} $\gets$ \code{None}
\EndIf
\State \code{noise_analysis} $\gets$ \code{quantify_effect_snr(src, src_noisy, ref, ref_noisy)}
\State \Return \code{(src_noisy, noise_analysis)}

\end{algorithmic}
\end{algorithm}

\begin{algorithm}[H]
\caption{\code{quantify_effect_snr(...)}}
\label{def:quantify-noise-utility}
\textbf{Purpose.} Computes the effect-to-noise signal-to-noise ratio between a clean trajectory, its noisy version, and optionally a matched reference trajectory.

\textbf{Used by.} Noise analysis during materialization and cheap-filter metrics.
\par\smallskip
\begin{algorithmic}[1]
\State \code{snr_analysis} $\gets$ \code{[]}
\ForAll{observed non-time signals \code{signal}}
    \If{\code{ref} is not \code{None}}
        \State \code{clean_effect_signal} $\gets$ \code{src[signal]} $-$ \code{ref[signal]}
        \If{\code{ref_noisy} is not \code{None}}
            \State \code{src_noise_signal} $\gets$ \code{src_noisy[signal]} $-$ \code{src[signal]}
            \State \code{ref_noise_signal} $\gets$ \code{ref_noisy[signal]} $-$ \code{ref[signal]}
            \State \code{noise_signal} $\gets$ \code{src_noise_signal} $-$ \code{ref_noise_signal}
        \Else
            \State \code{noise_signal} $\gets$ \code{src_noisy[signal]} $-$ \code{src[signal]}
        \EndIf
    \Else
        \State \code{clean_effect_signal} $\gets$ \code{src[signal]}
        \State \code{noise_signal} $\gets$ \code{src_noisy[signal]} $-$ \code{src[signal]}
    \EndIf
    \State append \code{SNR_db(clean_effect_signal, noise_signal)} to \code{snr_analysis}
\EndFor
\State \Return \code{snr_analysis}
\end{algorithmic}
\end{algorithm}

Where 
\begin{DocCode}
    SNR_db(effect_signal, noise) = 20 log10(RMS(effect_signal) / RMS(noise))
\end{DocCode}
\begin{itemize}
    \item \code{effect_signal} is the intervention effect relative to the reference trajectory, not the raw signal amplitude.
    \item Signals undergo preprocessing before the formula is applied, following \code{docs/overleaf_paper/src/appendix/appendix_environment.tex::*::Signal preprocessing}.
    \item \code{quantify_effect_snr} also accepts \textbf{list} inputs for \code{src}, \code{src_noisy}, \code{ref=None}, and \code{ref_noisy=None}.
    \item For list inputs, it returns the mean SNR by first averaging the linear effect-to-noise ratios over the list.
    \item The returned SNR in dB is then computed as \code{20*log10(mean effect-to-noise ratio)}.
    \end{itemize}
\end{document}
